% GENERATED FILE. Edit canonical lessons, then run npm run generate:books.
% canonical-source-hash: fnv1a64:386a641cbd43d968
% canonical-lessons: MR-C321-pangat, MR-C321-pahuncar, MR-C321-snehsammelan, MR-C321-upasthit-rahne, MR-C321-hajeri

\chapter{Guests and a Feast}
\label{ch:mr-guests-and-a-feast}
\hlchaptermodality{321}

\begin{chapteropening}
\textbf{By the end of this chapter:} \emph{I can talk about guests, attendance and a shared meal.}

Complete the last lesson of chapter 321: I can talk about guests, attendance and a shared meal.
\end{chapteropening}

\section[\texorpdfstring{\mr{पंगत} (paṅgat)}{पंगत (paṅgat)}]{\mr{पंगत} (paṅgat) — a row of diners}

\label{lesson:MR-C321-pangat}



\begin{quote}
Before the new one: say the Marathi for a pilgrimage, then the Marathi for an auspicious time.
\end{quote}

\subsection*{You'll want to know: \mr{पंगत}}
\textbf{\mr{पंगत}} — \emph{paṅgat} — \textquotedblleft{}a row of diners\textquotedblright{}.

\begin{quote}
\mr{लग्नात} \mr{पाहुण्यांच्या} \mr{पंगती} \mr{बसल्या}.
\end{quote}

\begin{quote}
\emph{lagnāt pāhuṇyāñcyā paṅgtī basalyā.}
\end{quote}

\begin{quote}
At the wedding the guests sat down in rows to eat.
\end{quote}

Feminine. Plural: \mr{पंगती}, \emph{paṅgtī}.

\mr{पंगत} \mr{बसणे}, a row sits down to a meal; plural \mr{पंगती}.

\begin{cousinweb}
\textbf{\mr{पंगत}} (\emph{paṅgat}) is usually traced to Sanskrit \emph{paṅkti}, 'a row, a troop'. CDIAL lists the inherited Marathi \emph{pā̃t}, 'a row', under that word (CDIAL 7646) but not \emph{paṅgat} itself, so its path is not settled.
\end{cousinweb}

\begin{grammarlens}[title={What you've built}]
One more word for telling what happened.
\end{grammarlens}

\subsection*{Your turn}
\begin{itemize}
\raggedright
  \item \emph{Say it:} \emph{paṅgat}
  \item \emph{Say it:} \emph{paṅgat}, once more
  \item \emph{Say it:} \emph{lagnāt pāhuṇyāñcyā paṅgtī basalyā.} — At the wedding the guests sat down in rows to eat.
  \item \emph{Recall:} say the Marathi for a pilgrimage, then the Marathi for an auspicious time, then say \emph{paṅgat} again
\end{itemize}

\begin{tcolorbox}[breakable,colback=teal!4,colframe=teal!35!black,title={Before you move on}]
What does \mr{पंगत} mean? (\textquotedblleft{}A row of diners\textquotedblright{}.) Say it once more.
\end{tcolorbox}
\section[\texorpdfstring{\mr{पाहुणचार} (pāhuṇcār)}{पाहुणचार (pāhuṇcār)}]{\mr{पाहुणचार} (pāhuṇcār) — hosting guests; hospitality}

\label{lesson:MR-C321-pahuncar}



\begin{quote}
Before the new one: say the Marathi for an auspicious time, then the Marathi for a row of diners.
\end{quote}

\subsection*{You'll want to know: \mr{पाहुणचार}}
\textbf{\mr{पाहुणचार}} — \emph{pāhuṇcār} — \textquotedblleft{}hosting guests; hospitality\textquotedblright{}.

\begin{quote}
\mr{त्या} \mr{गावातल्या} \mr{लोकांचा} \mr{पाहुणचार} \mr{आम्ही} \mr{कधीच} \mr{विसरणार} \mr{नाही}.
\end{quote}

\begin{quote}
\emph{tyā gāvātalyā lokāñcā pāhuṇcār āmhī kadhīc visarṇār nāhī.}
\end{quote}

\begin{quote}
We will never forget the hospitality of the people of that village.
\end{quote}

Masculine.

Built on \mr{पाहुणा}, a guest.

\begin{cousinweb}
\textbf{\mr{पाहुणचार}} (\emph{pāhuṇcār}) is built on \emph{pāhuṇā}, 'a guest', inherited from Sanskrit \emph{prāhuṇa} (CDIAL 8973), with the ending \emph{-cār}.
\end{cousinweb}

\begin{grammarlens}[title={What you've built}]
One more word for telling what happened.
\end{grammarlens}

\subsection*{Your turn}
\begin{itemize}
\raggedright
  \item \emph{Say it:} \emph{pāhuṇcār}
  \item \emph{Say it:} \emph{pāhuṇcār}, once more
  \item \emph{Say it:} \emph{tyā gāvātalyā lokāñcā pāhuṇcār āmhī kadhīc visarṇār nāhī.} — We will never forget the hospitality of the people of that village.
  \item \emph{Recall:} say the Marathi for an auspicious time, then the Marathi for a row of diners, then say \emph{pāhuṇcār} again
\end{itemize}

\begin{tcolorbox}[breakable,colback=teal!4,colframe=teal!35!black,title={Before you move on}]
What does \mr{पाहुणचार} mean? (\textquotedblleft{}Hosting guests; hospitality\textquotedblright{}.) Say it once more.
\end{tcolorbox}
\section[\texorpdfstring{\mr{स्नेहसंमेलन} (snehsammelan)}{स्नेहसंमेलन (snehsammelan)}]{\mr{स्नेहसंमेलन} (snehsammelan) — a get-together; an annual social}

\label{lesson:MR-C321-snehsammelan}



\begin{quote}
Before the new one: say the Marathi for a row of diners, then the Marathi for hosting guests.
\end{quote}

\subsection*{You'll want to know: \mr{स्नेहसंमेलन}}
\textbf{\mr{स्नेहसंमेलन}} — \emph{snehsammelan} — \textquotedblleft{}a get-together; an annual social\textquotedblright{}.

\begin{quote}
\mr{शाळेचं} \mr{वार्षिक} \mr{स्नेहसंमेलन} \mr{पुढच्या} \mr{आठवड्यात} \mr{आहे}.
\end{quote}

\begin{quote}
\emph{śāḷecaṁ vārṣik snehsammelan puḍhacyā āṭhvaḍyāt āhe.}
\end{quote}

\begin{quote}
The school's annual get-together is next week.
\end{quote}

Neuter.

\mr{स्नेह}, affection, + \mr{संमेलन}, a gathering.

\begin{cousinweb}
\textbf{\mr{स्नेहसंमेलन}} (\emph{snehsammelan}) is a modern compound built from Sanskrit \emph{sneha}, 'affection', and \emph{sammelana}, 'a meeting, a gathering' (Monier-Williams).
\end{cousinweb}

\begin{grammarlens}[title={What you've built}]
One more word for telling what happened.
\end{grammarlens}

\subsection*{Your turn}
\begin{itemize}
\raggedright
  \item \emph{Say it:} \emph{snehsammelan}
  \item \emph{Say it:} \emph{snehsammelan}, once more
  \item \emph{Say it:} \emph{śāḷecaṁ vārṣik snehsammelan puḍhacyā āṭhvaḍyāt āhe.} — The school's annual get-together is next week.
  \item \emph{Recall:} say the Marathi for a row of diners, then the Marathi for hosting guests, then say \emph{snehsammelan} again
\end{itemize}

\begin{tcolorbox}[breakable,colback=teal!4,colframe=teal!35!black,title={Before you move on}]
What does \mr{स्नेहसंमेलन} mean? (\textquotedblleft{}A get-together; an annual social\textquotedblright{}.) Say it once more.
\end{tcolorbox}
\section[\texorpdfstring{\mr{उपस्थित} \mr{राहणे} (upasthit rāhṇe)}{उपस्थित राहणे (upasthit rāhṇe)}]{\mr{उपस्थित} \mr{राहणे} (upasthit rāhṇe) — to attend; to be present}

\label{lesson:MR-C321-upasthit-rahne}



\begin{quote}
Before the new one: say the Marathi for hosting guests, then the Marathi for a get-together.
\end{quote}

\subsection*{You'll want to know: \mr{उपस्थित} \mr{राहणे}}
\textbf{\mr{उपस्थित} \mr{राहणे}} — \emph{upasthit rāhṇe} — \textquotedblleft{}to attend; to be present\textquotedblright{}.

\begin{quote}
\mr{कार्यक्रमाला} \mr{गावातले} \mr{सगळे} \mr{लोक} \mr{उपस्थित} \mr{राहिले}.
\end{quote}

\begin{quote}
\emph{kāryakramālā gāvātle sagḷe lok upasthit rāhile.}
\end{quote}

\begin{quote}
All the people of the village attended the event.
\end{quote}

Past (he …): \mr{तो} \mr{उपस्थित} \mr{राहिला}, \emph{to upasthit rāhilā}.

The event takes -\mr{ला}: \mr{सभेला} \mr{उपस्थित} \mr{राहणे}.

\begin{cousinweb}
\textbf{\mr{उपस्थित} \mr{राहणे}} (\emph{upasthit rāhṇe}) joins \emph{upasthit}, a tatsama taken from Sanskrit \emph{upasthita}, 'present, near at hand' (Monier-Williams), to \emph{rāhṇe}, 'to stay, to be'.
\end{cousinweb}

\begin{grammarlens}[title={What you've built}]
One more word for telling what happened.
\end{grammarlens}

\subsection*{Your turn}
\begin{itemize}
\raggedright
  \item \emph{Say it:} \emph{upasthit rāhṇe}
  \item \emph{Say it:} \emph{upasthit rāhṇe}, once more
  \item \emph{Say it:} \emph{kāryakramālā gāvātle sagḷe lok upasthit rāhile.} — All the people of the village attended the event.
  \item \emph{Recall:} say the Marathi for hosting guests, then the Marathi for a get-together, then say \emph{upasthit rāhṇe} again
\end{itemize}

\begin{tcolorbox}[breakable,colback=teal!4,colframe=teal!35!black,title={Before you move on}]
What does \mr{उपस्थित} \mr{राहणे} mean? (\textquotedblleft{}To attend; to be present\textquotedblright{}.) Say it once more.
\end{tcolorbox}
\section[\texorpdfstring{\mr{हजेरी} (hajerī)}{हजेरी (hajerī)}]{\mr{हजेरी} (hajerī) — attendance; presence}

\label{lesson:MR-C321-hajeri}



\begin{quote}
Before the new one: say the Marathi for a get-together, then the Marathi for to attend.
\end{quote}

\subsection*{You'll want to know: \mr{हजेरी}}
\textbf{\mr{हजेरी}} — \emph{hajerī} — \textquotedblleft{}attendance; presence\textquotedblright{}.

\begin{quote}
\mr{गुरुजींनी} \mr{वर्गात} \mr{हजेरी} \mr{घेतली}.
\end{quote}

\begin{quote}
\emph{gurujīnnī vargāt hajerī ghetlī.}
\end{quote}

\begin{quote}
The teacher took attendance in class.
\end{quote}

Feminine.

\mr{हजेरी} \mr{घेणे}, to take attendance; \mr{हजेरी} \mr{लावणे}, to put in an appearance.

\begin{cousinweb}
\textbf{\mr{हजेरी}} (\emph{hajerī}) is a loan from Persian \emph{hāzirī}, 'attendance', usually traced to Arabic \emph{ḥāḍir}, 'present'.
\end{cousinweb}

\begin{grammarlens}[title={What you've built}]
One more word for telling what happened.
\end{grammarlens}

\subsection*{Your turn}
\begin{itemize}
\raggedright
  \item \emph{Say it:} \emph{hajerī}
  \item \emph{Say it:} \emph{hajerī}, once more
  \item \emph{Say it:} \emph{gurujīnnī vargāt hajerī ghetlī.} — The teacher took attendance in class.
  \item \emph{Recall:} say the Marathi for a get-together, then the Marathi for to attend, then say \emph{hajerī} again
\end{itemize}

\begin{tcolorbox}[breakable,colback=teal!4,colframe=teal!35!black,title={Before you move on}]
What does \mr{हजेरी} mean? (\textquotedblleft{}Attendance; presence\textquotedblright{}.) Say it once more.
\end{tcolorbox}
